\chapter{Conclusion}

Understanding cyclist interactions during complex maneuvers is crucial for improving urban traffic safety, infrastructure design, and modeling of behavioral dynamics. While prior research has focused on isolated kinematic measures or short-term interactions, this thesis addresses the gap by capturing structured, recurring patterns in maneuver execution, linking micro-level motion dynamics to broader behavioral organization.

This work developed a hierarchical, unsupervised framework for identifying cyclist riding styles during complex maneuvers, such as overtaking and following interactions. The framework operates in three stages. Stage one decomposes each maneuver into short, overlapping temporal windows, capturing local riding dynamics through kinematic, traffic, and infrastructure features. This stage identifies micro-level behavioral regimes, revealing patterns such as stable cruising versus reactive, volatile control. Stage two aggregates these regime sequences into maneuver-level representations, summarizing temporal persistence, regime proportions, entropy, and scenario exposure. This aggregation preserves both the diversity and temporal organization of micro-regimes, enabling the subsequent clustering step to identify coherent execution styles. Stage three evaluates the relationship between these execution styles and environmental context, examining traffic exposure and infrastructure features descriptively to explore conditions under which different behaviors occur.

The framework revealed clear distinctions between overtaking and following behavior. Overtaking maneuvers were characterized by two micro-regimes—stable and volatile—and generally behaved as event-driven sequences, dominated by persistent, cohesive motion once initiated. Five distinct overtaking execution styles emerged, differentiating speed, acceleration strategy, lateral dynamics, and scenario context. Intersection maneuvers exhibited increased variability, particularly for accelerating or cautious overtakes, while free-flow overtakes tended to be smoother and more consistent. Following maneuvers, in contrast, displayed gradual variations in longitudinal control, with micro-regimes reflecting continuous adjustments rather than discrete states. Six execution styles captured differences in gap size, responsiveness, duration, and speed adaptation. Temporal aggregation highlighted that following maneuvers involve frequent short-term adjustments, whereas overtaking maneuvers typically follow long persistent runs, indicating fundamentally different temporal organizations between maneuver types.

Analysis of environmental context indicated that infrastructure and traffic exposure contributed only modestly to execution style differentiation. Overtaking styles showed minor differences in pedestrian or vehicle interaction at intersections, while following styles appeared primarily governed by cyclist interaction dynamics rather than contextual factors. This reinforces the insight that the internal dynamics of maneuvers play a central role in shaping execution styles.

The thesis contributes methodologically and empirically. Methodologically, it introduces a modular, multi-stage framework that transforms raw trajectory interactions into meaningful behavioral representations at both micro and maneuver levels. Empirically, it identifies distinct execution styles for overtaking and following, offering a nuanced understanding of internal kinematic patterns and temporal structures that were previously uncharacterized. These contributions provide a foundation for more informed traffic modeling, simulation of cyclist behavior, and urban safety interventions.

Limitations include the choice of window segmentation strategy, which may overlook subtle transitions between behavioral states, and some features, such as rotation fluctuations, were not corrected for geometric effects, potentially exaggerating variability in complex environments. Certain maneuver clusters—particularly sharp-steering overtakes—had small sample sizes, limiting interpretability, while following clusters showed moderate separation, reflecting continuous rather than discrete execution styles. Additionally, the descriptive analysis of environmental context, though informative, does not establish causal relationships.

Future work should explore the generalizability of the framework across datasets and urban contexts, assess alternative feature sets and segmentation strategies, and extend the approach to additional maneuvers such as turning, merging, or post-signal accelerations. Phase-specific analyses and modeling of decision processes could further examine why cyclists adopt particular execution styles under different interaction scenarios.


In summary, this thesis provides a comprehensive framework for identifying cyclist execution styles from trajectory data, linking micro-level dynamics to maneuver-level behavior. The results offer both methodological and empirical insights, advancing our understanding of cyclist interactions and providing a foundation for applications in traffic modeling, safety assessment, and urban mobility planning.

